% Generated by book/tools/notebook_to_tex.py. Do not edit by hand.

\chapter[이진 분류 (Binary Classification \& BCE)]{이진 분류 (Binary Classification \& BCE)}

\label{ch:03}

\markboth{이진 분류 (Binary Classification \& BCE)}{이진 분류 (Binary Classification \& BCE)}

\chaptermeta{03_sklearn_binary/03_sklearn_binary.ipynb}{https://colab.research.google.com/github/yoon-gu/neuqes-101/blob/master/03_sklearn_binary/03_sklearn_binary.ipynb}{logit, sigmoid, BCE가 만나는 방식}



\index{1Binary classification@Binary classification}

\index{1BCEWithLogitsLoss@BCEWithLogitsLoss}

\index{1sigmoid@sigmoid}

\index{1LogisticRegression@LogisticRegression}

\index{1predict proba@predict\_proba}

\index{0이진 분류@이진 분류}

\index{0로짓@로짓}

\index{0시그모이드@시그모이드}

\index{0예측 확률@예측 확률}

\index{0임계값@임계값}

\index{0정밀도@정밀도}

\index{0재현율@재현율}

\index{0혼동 행렬@혼동 행렬}

\index{1binary cross entropy@binary cross entropy}

\index{1log loss@log loss}

\index{1probability threshold@probability threshold}

\index{1threshold tuning@threshold tuning}

\index{1precision score@precision\_score}

\index{1recall score@recall\_score}

\index{1f1 score@f1\_score}

\index{1accuracy score@accuracy\_score}

\index{1class weight@class\_weight}

\index{1decision boundary@decision boundary}

\index{1positive class@positive class}

\index{1negative class@negative class}

\index{0이진 교차 엔트로피@이진 교차 엔트로피}

\index{0로그 손실@로그 손실}

\index{0확률 임계값@확률 임계값}

\index{0임계값 튜닝@임계값 튜닝}

\index{0양성 클래스@양성 클래스}

\index{0음성 클래스@음성 클래스}

\index{0결정 경계@결정 경계}

\index{0정확도@정확도}

\index{0F1 점수@F1 점수}



\section{변화추적표}

\begin{booktable}{3장 변화추적표}{tab:ch03-01}
\begin{adjustbox}{max width=\textwidth}
\begin{tabular}{@{}lllllll@{}}
\toprule
Ch & 모델 & 토크나이저 & 데이터 & Output Head & Activation & Loss \\
\midrule

1 & (TF-IDF) & \inlinecode{TfidfVectorizer()} & Yelp 5,000 샘플 & --- & --- & --- \\
2 & \inlinecode{LinearRegression()} & \inlinecode{TfidfVectorizer()} & Yelp (별점 1-5) & (1차원) & 없음 & \inlinecode{MSELoss} \\
\textbf{3 \(\leftarrow\) 여기} & \inlinecode{LogisticRegression()} & \inlinecode{TfidfVectorizer()} & Yelp 이진화 (4-5\(\to\)1, 1-2\(\to\)0, 3 제외) & (1차원) & \textbf{sigmoid} & \textbf{\inlinecode{BCEWithLogitsLoss}} (sklearn: log loss) \\
\bottomrule
\end{tabular}
\end{adjustbox}
\end{booktable}

전체 \ref{ch:20}장 표는 \href{https://github.com/yoon-gu/neuqes-101\#챕터별-변화추적표}{루트 README.md}를 참고하세요.



\section{변경점: \ref{ch:02}장 대비}

\begin{booktable}{3장 변경점 요약}{tab:ch03-02}
\begin{adjustbox}{max width=\textwidth}
\begin{tabular}{@{}lll@{}}
\toprule
축 & \ref{ch:02}장 & \ref{ch:03}장 \\
\midrule

모델 & \inlinecode{LinearRegression} & \textbf{\inlinecode{LogisticRegression}} (출력에 sigmoid 내장) \\
Activation & 없음 & \textbf{sigmoid} \\
Loss & \inlinecode{MSELoss} & \textbf{\inlinecode{BCEWithLogitsLoss}} \\
라벨 & float (1-5) & \textbf{int (0/1)} \\
데이터 & Yelp 별점 1-5 & \textbf{Yelp 이진화} (3 제외) \\
토크나이저 & TF-IDF & TF-IDF (그대로) \\
\bottomrule
\end{tabular}
\end{adjustbox}
\end{booktable}

겉보기엔 다섯 곳이 바뀌었지만, 한 가지 결정 --- \textbf{“회귀 \(\to\) 이진 분류”} --- 이 자동으로 끌어오는 결과입니다. 분류 패러다임은 (출력 형태 = sigmoid, loss = BCE, 라벨 = 0/1, 데이터 = 이진화)가 한 묶음으로 따라옵니다.

\textbf{왜 이렇게 묶이나?} 회귀에서 본 한계 --- “출력이 {[}0, 1{]}을 못 지킨다” --- 를 모델 단계에서 강제로 해결하는 게 sigmoid입니다. 그러면 출력은 자연스럽게 “1일 확률”로 해석되고, 그 확률에 어울리는 loss가 BCE이고, 라벨도 0/1 정수로 단순해집니다.



\section{\texorpdfstring{손실 함수의 변화 --- 이진 교차 엔트로피 등장}{손실 함수의 변화 --- 이진 교차 엔트로피 등장}}

\textbf{Binary Cross-Entropy} 는 모델이 출력한 확률 \(\hat p_i\) 와 정답 \(y_i \in \{0, 1\}\) 사이의 차이를 잽니다.

\begin{equation}
\label{eq:ch03-01}
L = -\frac{1}{N}\sum_{i=1}^{N}\left[\,y_i \log \hat p_i + (1 - y_i)\log(1 - \hat p_i)\,\right]
\end{equation}

식~\eqref{eq:ch03-01}은 정답 라벨과 예측 확률 사이의 Binary Cross-Entropy를 정의합니다.

핵심: \(y_i = 1\)이면 첫 항 \(-\log \hat p_i\)만 살고, \(y_i = 0\)이면 둘째 항 \(-\log(1 - \hat p_i)\)만 살아남습니다. 즉 한 샘플당 항상 한 항만 작동합니다.

\textbf{숫자로 감 잡기} (정답이 \(y = 1\)인 샘플 한 개, \(N = 1\)):

\begin{booktable}{3장 손실 함수의 변화 - 이진 교차 엔트로피 등장 표}{tab:ch03-03}
\begin{adjustbox}{max width=\textwidth}
\begin{tabular}{@{}lll@{}}
\toprule
정답 \(y\) & 예측 확률 \(\hat p\) & 손실 \(-\log \hat p\) \\
\midrule

1 & 0.9 & 0.105 \\
1 & 0.5 & 0.693 \\
1 & 0.1 & \textbf{2.303} \\
1 & 0.01 & \textbf{4.605} \\
\bottomrule
\end{tabular}
\end{adjustbox}
\end{booktable}

확률이 정답에서 멀어질수록 --- 특히 0에 가까워질수록 --- 손실이 \textbf{로그 스케일로 폭증} 합니다. 자신 있게 틀린 예측을 강하게 처벌하는 게 BCE의 성격입니다.

\textbf{\inlinecode{BCEWithLogits}의 “Logits” 의미}: 모델 마지막 단의 raw 점수(logit) \(z = w^\top x + b\)를 sigmoid에 넣기 \emph{전} 의 값을 의미합니다. PyTorch의 \inlinecode{BCEWithLogitsLoss}는 logit을 받아 내부에서 sigmoid + BCE를 한 번에 계산하는데, 따로 sigmoid를 통과시킨 뒤 BCE를 적용하는 것보다 수치적으로 안정적입니다.

\Needspace{8\baselineskip}
\begin{lstlisting}
criterion = nn.BCEWithLogitsLoss()
loss = criterion(logits, targets.float())

from sklearn.linear_model import LogisticRegression
model = LogisticRegression(max_iter=1000)
model.fit(X, y)
\end{lstlisting}



\section{토크나이저 노트}

이 장의 토크나이저는 \textbf{\ref{ch:01}·\ref{ch:02}장과 동일한 \inlinecode{TfidfVectorizer}} 입니다. 입력 표현은 그대로고, 변화는 모델 출력단·loss·라벨에서만 일어납니다.

\begin{previewBox}{미리보기}
\textbf{다음 장(\ref{ch:04}장)}: 같은 TF-IDF 그대로. 변하는 건 출력이 5차원으로 늘어나고 sigmoid가 softmax로 바뀌는 것뿐.
\end{previewBox}



\Needspace{5\baselineskip}
\begin{lstlisting}
!pip install -q datasets scikit-learn pandas matplotlib
\end{lstlisting}

\begin{codeRead}
\textbf{1행}의 \inlinecode{!pip install -q datasets...}에서는 Colab 실행에 필요한 패키지를 설치합니다.
\end{codeRead}



\Needspace{11\baselineskip}
\begin{lstlisting}
import numpy as np
import pandas as pd
import matplotlib.pyplot as plt
from datasets import load_dataset
from sklearn.feature_extraction.text import TfidfVectorizer
from sklearn.linear_model import LogisticRegression
from sklearn.model_selection import train_test_split
from sklearn.metrics import (
    accuracy_score, classification_report, confusion_matrix,
    log_loss, precision_recall_fscore_support,
)

plt.rcParams["axes.unicode_minus"] = False

\end{lstlisting}

\noindent\emph{앞 블록은 필요한 라이브러리와 기본 설정을 준비하는 단계입니다. 이어지는 블록에서는 데이터를 불러오고 학습에 맞는 형태로 정리하는 단계로 넘어갑니다.}

\Needspace{9\baselineskip}
\begin{lstlisting}[firstnumber=15]
dataset = load_dataset("Yelp/yelp_review_full")
SAMPLE_SIZE = 5000
ds = dataset["train"].shuffle(seed=42).select(range(SAMPLE_SIZE))
df = ds.to_pandas()
df["star"] = df["label"] + 1   # 0-4 -> 1-5
print(f"Total samples: {len(df)}")
print(df["star"].value_counts().sort_index())
\end{lstlisting}

\begin{codeRead}
\inlinecode{numpy, pandas, matplotlib, datasets, scikit-learn} 같은 주요 패키지를 불러옵니다.
\end{codeRead}

\noindent\textbf{출력.}
\begin{lstlisting}[style=bookoutput]
If the error persists, please let us know by opening an issue on GitHub (ht...
Total samples: 5000
star
1    1017
2    1027
3     960
4    1021
5     975
Name: count, dtype: int64
\end{lstlisting}
\par\vspace{0.9em}



\Needspace{11\baselineskip}
\begin{lstlisting}
df_bin = df[df["star"] != 3].copy()
df_bin["y"] = (df_bin["star"] >= 4).astype(int)

print(
    f"Binarized samples: "
    f"{len(df_bin)}  (3-star excluded)"
)
print(
    f"Class distribution:\n"
    f"{df_bin['y'].value_counts().sort_index()}"
)
print(f"Positive rate: {df_bin['y'].mean():.1%}")
\end{lstlisting}

\begin{codeRead}
\textbf{1행}의 \inlinecode{df\_bin = ...}에서는 이후 분석에서 사용할 값을 준비합니다. \textbf{2행}의 \inlinecode{df\_bin["y"] = (df\_bin["st...}에서는 이후 분석에서 사용할 값을 준비합니다.
\end{codeRead}

\noindent\textbf{출력.}
\begin{lstlisting}[style=bookoutput]
Binarized samples: 4040  (3-star excluded)
Class distribution:
y
0    2044
1    1996
Name: count, dtype: int64
Positive rate: 49.4%
\end{lstlisting}
\par\vspace{0.9em}



\Needspace{11\baselineskip}
\begin{lstlisting}
X_text_train, X_text_test, y_train, y_test = train_test_split(
    df_bin["text"], df_bin["y"],
    test_size=0.2, random_state=42, stratify=df_bin["y"],
)

tfidf = TfidfVectorizer(max_features=10000)
X_train = tfidf.fit_transform(X_text_train)
X_test = tfidf.transform(X_text_test)

print(
    f"X_train: {X_train.shape}, y_train: "
    f"{y_train.shape}"
)
print(f"X_test:  {X_test.shape}, y_test:  {y_test.shape}")
\end{lstlisting}

\begin{codeRead}
\textbf{1--4행}의 \inlinecode{X\_text\_train, X\_text\_test...}에서는 학습용 데이터와 평가용 데이터를 분리합니다. \textbf{6행}의 \inlinecode{tfidf = ...}에서는 TF-IDF 기반 벡터라이저를 만듭니다. \textbf{7행}의 \inlinecode{X\_train = ...}에서는 텍스트나 라벨을 모델이 처리할 수 있는 수치 표현으로 변환합니다. \textbf{8행}의 \inlinecode{X\_test = ...}에서는 텍스트나 라벨을 모델이 처리할 수 있는 수치 표현으로 변환합니다.
\end{codeRead}

\noindent\textbf{출력.}
\begin{lstlisting}[style=bookoutput]
X_train: (3232, 10000), y_train: (3232,)
X_test:  (808, 10000), y_test:  (808,)
\end{lstlisting}
\par\vspace{0.9em}



\section{\texorpdfstring{실습: \inlinecode{LogisticRegression}으로 이진 분류}{실습: LogisticRegression으로 이진 분류}}

\inlinecode{LogisticRegression}이 내부에서 하는 일은 두 단계입니다.

\begin{enumerate}
\def\labelenumi{\arabic{enumi}.}
\tightlist
\item
  logit 계산: \(z = w^\top x + b\) (\ref{ch:02}장과 똑같은 선형 결합)
\item
  sigmoid로 확률 변환: \(\hat p = \sigma(z) = \dfrac{1}{1 + e^{-z}}\)
\end{enumerate}

학습은 BCE를 최소화하는 \(w, b\)를 찾는 것이고, 예측은 \(\hat p \geq 0.5\)를 기준으로 0/1로 자릅니다.



\Needspace{10\baselineskip}
\begin{lstlisting}
model = LogisticRegression(max_iter=1000)
model.fit(X_train, y_train)

y_pred = model.predict(X_test)
print(
    f"Test accuracy: "
    f"{accuracy_score(y_test, y_pred):.4f}"
)
\end{lstlisting}

\begin{codeRead}
\textbf{1행}의 \inlinecode{model = ...}에서는 이번 실습에서 관찰할 모델 객체를 정의합니다. \textbf{2행}의 \inlinecode{model.fit(...)}에서는 학습 데이터로 모델 또는 변환기를 적합합니다. \textbf{4행}의 \inlinecode{y\_pred = ...}에서는 학습된 모델로 최종 예측값을 만듭니다.
\end{codeRead}

\noindent\textbf{출력.}
\begin{lstlisting}[style=bookoutput]
Test accuracy: 0.8639
\end{lstlisting}
\par\vspace{0.9em}



\Needspace{11\baselineskip}
\begin{lstlisting}
y_proba = model.predict_proba(X_test)
print(
    f"y_proba shape: "
    f"{y_proba.shape}  (per sample: [P(0), P(1)])"
)
print(f"\nFirst 5 predicted probabilities:")
print(
    pd.DataFrame(y_proba[:5], columns=["P(neg)", "P(pos)"]).round(4)
)

print(
    f"\nRow sums (should be 1): "
    f"{y_proba.sum(axis=1)[:5]}"
)
\end{lstlisting}

\begin{codeRead}
\textbf{1행}의 \inlinecode{y\_proba = ...}에서는 클래스별 예측 확률을 계산합니다.
\end{codeRead}

\noindent\textbf{출력.}
\begin{lstlisting}[style=bookoutput]
y_proba shape: (808, 2)  (per sample: [P(0), P(1)])

First 5 predicted probabilities:
   P(neg)  P(pos)
0  0.5523  0.4477
1  0.8999  0.1001
2  0.8409  0.1591
3  0.3347  0.6653
4  0.2420  0.7580

Row sums (should be 1): [1. 1. 1. 1. 1.]
\end{lstlisting}
\par\vspace{0.9em}



\section{해부: sigmoid는 logit을 어떻게 확률로 바꾸나}

\inlinecode{LogisticRegression}의 \inlinecode{coef\_}와 \inlinecode{intercept\_}는 \ref{ch:02}장 \inlinecode{LinearRegression}과 동일한 선형 가중치입니다. 차이는 그 위에 sigmoid가 한 번 더 붙는다는 것뿐.

직접 재현해 봅시다 --- sklearn의 \inlinecode{predict\_proba}가 정말 sigmoid에 logit을 넣은 결과인지 확인합니다.



\Needspace{11\baselineskip}
\begin{lstlisting}
logits = X_test @ model.coef_.T + model.intercept_   # shape: (
    N,
    1,
)
logits = logits.flatten()

proba_manual = 1 / (1 + np.exp(-logits))
proba_sklearn = y_proba[:, 1]   # P(y=1)

diff = np.abs(proba_manual - proba_sklearn).max()
print(f"Max diff (manual vs sklearn): {diff:.2e}")

print(f"\nManual first 5:  {proba_manual[:5].round(4)}")
print(f"sklearn first 5: {proba_sklearn[:5].round(4)}")
\end{lstlisting}

\begin{codeRead}
\textbf{1--4행}의 \inlinecode{logits = ...}에서는 이후 분석에서 사용할 값을 준비합니다. \textbf{5행}의 \inlinecode{logits = ...}에서는 이후 분석에서 사용할 값을 준비합니다. \textbf{7행}의 \inlinecode{proba\_manual = ...}에서는 이후 분석에서 사용할 값을 준비합니다. \textbf{8행}의 \inlinecode{proba\_sklearn = ...}에서는 이후 분석에서 사용할 값을 준비합니다. \textbf{10행}의 \inlinecode{diff = ...}에서는 두 계산 결과가 얼마나 다른지 확인할 값을 만듭니다.
\end{codeRead}

\noindent\textbf{출력.}
\begin{lstlisting}[style=bookoutput]
Max diff (manual vs sklearn): 1.11e-16

Manual first 5:  [0.4477 0.1001 0.1591 0.6653 0.758 ]
sklearn first 5: [0.4477 0.1001 0.1591 0.6653 0.758 ]
\end{lstlisting}
\par\vspace{0.9em}



\Needspace{10\baselineskip}
\begin{lstlisting}
y_test_arr = y_test.values
p = proba_sklearn
manual_bce = -(y_test_arr * np.log(p) + (1 - y_test_arr) * np.log(1 - p)).mean()
sklearn_bce = log_loss(y_test, y_proba)

print(f"Manual BCE:  {manual_bce:.6f}")
print(f"sklearn BCE: {sklearn_bce:.6f}")
print(f"Diff:        {abs(manual_bce - sklearn_bce):.2e}")
\end{lstlisting}

\begin{codeRead}
\textbf{1행}의 \inlinecode{y\_test\_arr = ...}에서는 계산에 바로 쓸 수 있도록 배열 형태로 변환합니다. \textbf{2행}의 \inlinecode{p = ...}에서는 이후 분석에서 사용할 값을 준비합니다. \textbf{3행}의 \inlinecode{manual\_bce = ...}에서는 두 계산 결과가 얼마나 다른지 확인할 값을 만듭니다. \textbf{4행}의 \inlinecode{sklearn\_bce = ...}에서는 이후 분석에서 사용할 값을 준비합니다.
\end{codeRead}

\noindent\textbf{출력.}
\begin{lstlisting}[style=bookoutput]
Manual BCE:  0.383569
sklearn BCE: 0.383569
Diff:        0.00e+00
\end{lstlisting}
\par\vspace{0.9em}



\Needspace{11\baselineskip}
\begin{lstlisting}
print(
    classification_report(y_test, y_pred, target_names=["negative", "positive"])
)

cm = confusion_matrix(y_test, y_pred)
print(f"\nconfusion matrix:\n{cm}")
print(
    f"  TN={cm[0,0]}, FP={cm[0,1]}\n  FN={cm[1,0]}, TP="
    f"{cm[1,1]}"
)
\end{lstlisting}

\begin{codeRead}
\textbf{5행}의 \inlinecode{cm = ...}에서는 이후 분석에서 사용할 값을 준비합니다.
\end{codeRead}

\noindent\textbf{출력.}
\begin{lstlisting}[style=bookoutput]
precision    recall  f1-score   support

    negative       0.87      0.86      0.86       409
    positive       0.86      0.87      0.86       399

    accuracy                           0.86       808
   macro avg       0.86      0.86      0.86       808
weighted avg       0.86      0.86      0.86       808


confusion matrix:
[[352  57]
 [ 53 346]]
  TN=352, FP=57
  FN=53, TP=346
\end{lstlisting}
\par\vspace{0.9em}



\section{변형: 임계값(threshold)을 옮기면}

\inlinecode{predict()}는 기본적으로 \(\hat p \geq 0.5\)를 기준으로 0/1을 자릅니다. 이 임계값을 다른 값으로 옮기면 precision과 recall이 정반대로 움직입니다.

\begin{itemize}
\tightlist
\item
  임계값 ↑ (예: 0.7): “확실해야만 positive” \(\to\) \textbf{precision 상승}, recall 하락
\item
  임계값 \(\downarrow\) (예: 0.3): “조금만 의심돼도 positive” \(\to\) \textbf{recall 상승}, precision 하락
\end{itemize}

스팸 필터(precision 중시), 암 진단(recall 중시) 같은 도메인 요구에 따라 임계값을 조정합니다.



\Needspace{11\baselineskip}
\begin{lstlisting}
thresholds = [0.3, 0.4, 0.5, 0.6, 0.7]
proba_pos = y_proba[:, 1]

rows = []
for t in thresholds:
    y_pred_t = (proba_pos >= t).astype(int)
    p, r, f1, _ = precision_recall_fscore_support(
        y_test, y_pred_t, average="binary", zero_division=0
    )
    acc = accuracy_score(y_test, y_pred_t)
    rows.append({"threshold": t, "accuracy": acc, "precision": p, "recall": r, "f1": f1})

df_t = pd.DataFrame(rows).round(4)
print(df_t.to_string(index=False))
\end{lstlisting}

\noindent\textbf{출력.}
\begin{lstlisting}[style=bookoutput]
threshold  accuracy  precision  recall     f1
       0.3    0.7686     0.6840  0.9875 0.8082
       0.4    0.8527     0.7929  0.9499 0.8643
       0.5    0.8639     0.8586  0.8672 0.8628
       0.6    0.8366     0.9159  0.7368 0.8167
       0.7    0.7772     0.9700  0.5664 0.7152
\end{lstlisting}
\par\vspace{0.9em}



\section{이 장에 등장한 라이브러리}

\begin{booktable}{3장 새로 등장한 라이브러리}{tab:ch03-04}
\begin{adjustbox}{max width=\textwidth}
\begin{tabular}{@{}lll@{}}
\toprule
이름 & 한 줄 설명 & 다음 장에서 \\
\midrule

\inlinecode{sklearn.linear\_model.LogisticRegression} & sigmoid + BCE 내장 이진/다중 분류 & \ref{ch:04}장에서 multi-class로 확장 \\
\inlinecode{sklearn.metrics.classification\_report} & accuracy/precision/recall/F1 한 번에 & 분류 장마다 계속 사용 \\
\inlinecode{sklearn.metrics.confusion\_matrix} & 혼동 행렬 & 다중 분류에서도 활용 \\
\inlinecode{sklearn.metrics.log\_loss} & BCE 평가 & \ref{ch:09}장 이후 BERT binary에서도 \\
\inlinecode{sklearn.metrics.precision\_recall\_fscore\_support} & 임계값별 지표 추적 & --- \\
\bottomrule
\end{tabular}
\end{adjustbox}
\end{booktable}



\section{체크포인트 질문}

\begin{enumerate}
\def\labelenumi{\arabic{enumi}.}
\tightlist
\item
  \inlinecode{LogisticRegression}이 내부에서 하는 두 단계(logit 계산, sigmoid)를 식으로 적어보세요.
\item
  BCE에서 정답이 \(y = 1\)이면 어느 항이 살아남고, 그 항이 예측 확률에 따라 어떻게 변하나요?
\item
  \inlinecode{predict\_proba}의 출력 shape가 \((N, 2)\)인 이유는? 두 열의 합은 항상 얼마여야 하나요?
\item
  임계값을 0.5에서 0.7로 올리면 precision과 recall은 어떤 방향으로 움직이고, 그 이유는 무엇인가요?
\end{enumerate}



\section{FAQ}

\begin{faqBox}{Q1. (이론) 왜 binary classification에서 MSE를 쓰면 안 되나요?}

작동은 하지만 두 가지 문제가 있습니다.

\begin{enumerate}
\def\labelenumi{\arabic{enumi}.}
\item
  \textbf{sigmoid + MSE는 비볼록(non-convex)}: 손실면이 매끈한 그릇 모양이 아니라 평탄한 구간이 생깁니다. local minima에 빠질 수 있고 학습이 불안정합니다. 반면 sigmoid + BCE는 볼록(convex)이라 전역 최적해로 수렴이 보장됩니다.
\item
  \textbf{gradient 소실(saturation)}: sigmoid 양 극단(출력이 0이나 1 근처)에서 도함수가 거의 0이 됩니다. MSE의 gradient는 그 도함수에 곱해지므로 함께 0으로 중단될 수 있습니다. BCE는 sigmoid와 결합되었을 때 gradient가 단순히 \(\hat p - y\)로 떨어져 saturation이 사라집니다.
\end{enumerate}

참고로 통계학에선 0/1 라벨을 그대로 \inlinecode{LinearRegression}으로 푸는 \textbf{Linear Probability Model(LPM)} 이 있긴 합니다. 단순 분석엔 쓰지만 출력이 {[}0, 1{]}을 안 지키고 BCE보다 학습 안정성이 떨어져 ML에선 거의 안 씁니다.

\end{faqBox}

\begin{faqBox}{Q2. (이론) sigmoid 대신 다른 활성화 함수(tanh, softmax)를 쓰면 어떻게 되나요?}

\begin{itemize}
\tightlist
\item
  \textbf{tanh}: 출력 범위가 \([-1, 1]\)입니다. 라벨을 -1/+1로 매핑하고 hinge loss 등을 쓰면 SVM과 비슷한 모델이 됩니다. 수학적으로는 \(\tanh(z) = 2\sigma(2z) - 1\)이라 sigmoid의 단순 변환이지만 라벨 컨벤션이 다릅니다.
\item
  \textbf{softmax}: 다중 클래스용 일반화. binary에 softmax를 쓰려면 출력 차원을 2로 늘리고 라벨도 one-hot으로 바꿔야 합니다 \(\to\) 이게 정확히 \ref{ch:09}장에서 다룰 “방식 B”입니다 (방식 A=sigmoid 1차원, 방식 B=softmax 2차원이 수학적으로 동등).
\item
  \textbf{ReLU/identity}: 출력이 {[}0, 1{]} 보장이 안 됩니다 --- 음수도 1 초과도 가능 \(\to\) BCE의 \(\log\)가 정의되지 않습니다.
\end{itemize}

\end{faqBox}

\begin{faqBox}{Q3. (실무) 클래스 불균형이 있으면 어떻게 하나요?}

세 가지 접근이 있습니다.

\begin{enumerate}
\def\labelenumi{\arabic{enumi}.}
\tightlist
\item
  \textbf{\texttt{class\_weight=\textquotesingle{}balanced\textquotesingle{}}}: sklearn에 한 줄로 적용. 소수 클래스의 손실에 더 큰 가중치.
\end{enumerate}

\Needspace{6\baselineskip}
\begin{lstlisting}
model = LogisticRegression(
    class_weight="balanced",
    max_iter=1000,
)
\end{lstlisting}

\begin{enumerate}
\def\labelenumi{\arabic{enumi}.}
\setcounter{enumi}{1}
\item
  \textbf{임계값 조정}: 확률 분포가 한쪽으로 쏠려 있어 0.5 임계값이 의미 없을 때, 위 셀의 threshold sweep으로 F1이 최대인 점을 찾습니다.
\item
  \textbf{데이터 레벨 처리}: SMOTE(소수 클래스 합성), undersampling(다수 클래스 줄이기). \inlinecode{imbalanced-learn} 라이브러리.
\end{enumerate}

\end{faqBox}

\begin{faqBox}{Q4. (실무) 임계값을 0.5 외에 다른 값으로 바꾸려면?}

\inlinecode{predict\_proba}로 확률을 얻은 뒤 직접 자르면 됩니다.

\Needspace{5\baselineskip}
\begin{lstlisting}
proba_pos = model.predict_proba(X_test)[:, 1]
y_pred_custom = (proba_pos >= 0.3).astype(int)
\end{lstlisting}

최적 임계값은 보통 검증 데이터에서 F1이 최대가 되는 지점, 또는 ROC 곡선의 Youden\textquotesingle s J 통계량(\inlinecode{tpr\ -\ fpr}이 최대인 지점)으로 잡습니다.

\Needspace{6\baselineskip}
\begin{lstlisting}
from sklearn.metrics import roc_curve
fpr, tpr, thr = roc_curve(y_test, proba_pos)
best_thr = thr[(tpr - fpr).argmax()]
print(f"Youden's J 기준 최적 임계값: {best_thr:.3f}")
\end{lstlisting}

\end{faqBox}

\begin{faqBox}{Q5. (이론) accuracy, precision, recall, F1 중 뭘 봐야 하나요?}

데이터 분포와 도메인 비용에 따라 다릅니다.

\begin{itemize}
\tightlist
\item
  \textbf{accuracy} (전체 맞춘 비율): 클래스가 균형 잡혔을 때만 의미 있음. 95\% 음성/5\% 양성 데이터에서 모두 음성이라 찍어도 95\%.
\item
  \textbf{precision} (positive 예측 중 진짜 positive 비율): \textbf{오탐 비용} 이 클 때. 스팸 필터, 광고 추천.
\item
  \textbf{recall} (실제 positive 중 잡아낸 비율): \textbf{놓치는 비용} 이 클 때. 암 진단, 사기 탐지.
\item
  \textbf{F1} (precision·recall의 조화 평균): 둘 다 중요할 때의 균형 지표. 어느 한 쪽이 0이면 F1도 0.
\end{itemize}

이 장의 Yelp 이진화는 클래스 불균형이 크지 않아(긍정 \(\approx\) 60\%) accuracy도 의미 있지만, 실무에선 항상 classification\_report 전체를 보는 습관이 안전합니다.

\end{faqBox}

\begin{faqBox}{Q6. (실무) sklearn \inlinecode{LogisticRegression}은 정규화(L2)가 기본인데 끄려면?}

\inlinecode{penalty=None}(0.22 이상) 또는 \inlinecode{C} 값을 매우 크게 설정합니다.

\Needspace{5\baselineskip}
\begin{lstlisting}
LogisticRegression(penalty=None, max_iter=1000)
LogisticRegression(C=1e10, max_iter=1000)
\end{lstlisting}

기본값은 \inlinecode{C=1.0}(L2 정규화)이고, 텍스트 분류처럼 feature가 많은 경우 정규화가 있어야 일반화 성능이 안정적입니다. 실무에선 거의 끄지 않습니다.
\end{faqBox}



\section{오류 실험 (선택)}

다음 코드를 돌려보고 결과를 예측해보세요. 0/1 라벨에 \inlinecode{LinearRegression}을 그대로 적용하면 무슨 일이 일어날까요?

\Needspace{11\baselineskip}
\begin{lstlisting}
from sklearn.linear_model import LinearRegression

lpm = LinearRegression()
lpm.fit(X_train, y_train)
pred_lpm = lpm.predict(X_test)

print(
    f"LPM 예측 범위: [{pred_lpm.min():.3f}, "
    f"{pred_lpm.max():.3f}]"
)
print(f"0 미만 예측: {(pred_lpm < 0).sum()}개")
print(f"1 초과 예측: {(pred_lpm > 1).sum()}개")

\end{lstlisting}

\noindent\emph{앞 블록은 필요한 라이브러리와 기본 설정을 준비하는 단계입니다. 이어지는 블록에서는 앞에서 만든 중간 값을 다음 계산으로 넘기는 단계로 넘어갑니다.}

\Needspace{8\baselineskip}
\begin{lstlisting}[firstnumber=14]
acc_lpm = accuracy_score(y_test, (pred_lpm >= 0.5).astype(int))
print(f"\nLPM accuracy (threshold 0.5): {acc_lpm:.4f}")
print(
    f"LogReg accuracy:               "
    f"{accuracy_score(y_test, y_pred):.4f}"
)
\end{lstlisting}

힌트: LPM은 출력이 {[}0, 1{]}을 안 지키지만, 임계값 0.5로 자르면 분류 성능 자체는 LogReg와 비슷할 수도 있습니다. 다만 “출력이 확률”이라는 해석을 잃습니다.



\begin{previewBox}{미리보기: 다음 장}

\textbf{\ref{ch:04}장. sigmoid와 softmax의 동등성 (Binary Classification) --- 같은 문제, 다른 표현}

\begin{itemize}
\tightlist
\item
  별점 1-5 를 5개 독립 클래스로 보고 \inlinecode{LogisticRegression()} 한 줄로 분류 (sklearn 이 multi-class 데이터에 자동으로 multinomial+softmax 적용)
\item
  출력이 1차원에서 \textbf{5차원} 으로 늘어나고, sigmoid 대신 \textbf{softmax} 가 붙음 (합 = 1 강제)
\item
  Loss 는 BCE 에서 \textbf{\inlinecode{CrossEntropyLoss}} (sklearn: multinomial log loss) 로 일반화
\item
  multinomial(softmax+CE) vs OvR(One-vs-Rest, K개 독립 binary) 비교 --- OvR 이 \ref{ch:06}장 multi-label 로 가는 다리
\end{itemize}
\end{previewBox}